\section*{Sources \& limitations}

本号では、公式発表・ドキュメント・モデルカード・リポジトリの記録・論文抄録・報道、コミュニティ上の反応を、それぞれの性質に合わせて読み分ける。公開情報が示す範囲と、まだ確かめられていない範囲を分けるためである。

\begin{tcolorbox}[enhanced,breakable,colback=SurveySoft,colframe=SurveyRule,boxrule=0.4pt,arc=1mm,left=3mm,right=3mm,top=2mm,bottom=2mm]
\textbf{公式・プロジェクト資料}\par
モデルやツールの提供条件、利用経路、機能、リリースの変更点を示す。提供元やプロジェクト自身が報告する速度・性能値は、独立測定の結果とは区別して読む。\par\smallskip
\textbf{論文抄録}\par
FUSEの枠組みと対照は抄録の範囲で受け止め、全文の手法やモデル一覧、点数、統計は未検証のまま残す。\par\smallskip
\textbf{コミュニティ上の反応}\par
W36の観測では、旗艦モデルのコンピューター利用やコード生成報告、公開モデルの手元環境での利用への関心が見られた。これは話題の広がりを示す背景情報であり、技術性能や互換性の証明には用いない。速報は文脈資料である。\par\smallskip
\textbf{まだ確かめる必要がある点}\par
重みの利用条件、プレビューと旗艦の境界、利用対象とセーフガード、ベンチマークの条件、買収の締結と規制、異なるモデル・環境への一般化などは、公開情報が増えるまで断定しない。
\end{tcolorbox}

本文では、各資料の主体、日付、提供経路、測定条件をできるだけ明示した。異なる主体が報告した数値や、条件の違うベンチマークを同じ尺度のランキングへ変換していない。Museの図表の数値は取り込まず、追加のモデルカード詳細は本号の事実確認には用いていない。

\nocite{astrasafety,astrapath,astralaunch,fuse,w36community,fermatscience,fermatgithub,k2horizon,glm53weights,fablemythos,gemini38flash,musespark,copilotharness,kilojetbrains,nvidiahf,atlas,pics,h3max}

\begin{claimboundary}[Temporal and claim boundary]
本文の基準時点は2026-09-04 18:00 America/New\_York。公式資料、リポジトリの記録、論文抄録、報道の伝える内容、コミュニティ上の観測は、それぞれ異なる範囲の主張を支える。提供元の速度値、報道の仕様、実験室の評価値、話題で見られた関心は相互に代替できず、測定条件が異なる値も単純比較できない。未解決の点は未解決のまま記載する。買収は合意・発表であり成立ではない。開放の言葉は誓いである。
\end{claimboundary}
